
\documentclass[12pt]{article}
\usepackage{amsmath}
\usepackage{amsfonts}
\usepackage{amssymb}
\usepackage{geometry}
\usepackage{graphicx}
\usepackage{accents}
\usepackage[thinc]{esdiff}
\usepackage{indentfirst}
\setlength{\parindent}{0pt}
\geometry{a4paper, margin=0.8in}
\pagenumbering{gobble}

\title{Homework assignment for 02YELMA \\ Chapter: Special theory of relativity}
\begin{document}
\date{}
\maketitle
\vspace{-0.8in}

\begin{enumerate}
    \item Consider a train moving at a speed of 50 km / h with respect to the ground and a kid inside the train throwing a ball at a speed of 10 km / h with respect to the train. We want to calculate the speed of the ball with respect to the ground. What is the percentage of error introduced when we calculate the speed of the ball with respect to the ground by simply adding the speeds of the train with respect to the ground and the ball with respect to the train?
    \item Muons are elementary particles similar to electrons but with a larger mass. Muons have a mean lifetime of about $2.2 \times 10^{-6}$ seconds in their rest frame. They are created in the upper atmosphere, approximately 15 kilometers above Earth's surface, because of cosmic-ray interactions.
    
    \textbf{Question:}
    Consider muons created at an altitude of 15 kilometers above the Earth's surface, traveling directly downward at a speed very close to the speed of light ($0.998c$).
    
    \begin{enumerate}
        \item \textbf{Without Special Relativity:}
        Calculate the distance a muon would travel before decaying if its lifetime is $2.2 \times 10^{-6}$ seconds. Based on this calculation, would the muons be able to reach the Earth's surface?
        
        \item \textbf{With Special Relativity:}
        Use the time dilation formula from special relativity to calculate the dilated lifetime of the muons as observed from the Earth's frame of reference. Show your work step by step.
        
        \begin{equation*}
            t = \frac{t'}{\sqrt{1 - \left(\frac{v}{c}\right)^2}}
        \end{equation*}
        
        where $t$ is the dilated lifetime (measured in the reference frame of the Earth), $t'$ is the lifetime in the reference frame of the muon ($2.2 \times 10^{-6}$ seconds), $v$ is the speed of the muon ($0.998c$) and $c$ is the speed of light.
        
        \item \textbf{Comparison and Conclusion:}
        Calculate the distance muons travel before decaying using the dilated lifetime. Compare this distance to the thickness of the atmosphere (15 kilometers). Based on this comparison, explain why special relativity is necessary to understand how muons created in the upper atmosphere can reach the Earth's surface.
    \end{enumerate}
    \item In the special theory of relativity, just as the coordinates of events and forces transform according to Lorentz transformations between relatively moving frames of reference, so too do the electric and magnetic fields. Accordingly, the fields $\vec{E}$ and $\vec{B}$ in the reference frame $S$ appear as follows in a frame $S'$ moving with constant velocity $v$ along the x-axis relative to $S$:
\[
\vec{E} = \begin{pmatrix}
E_x \\ E_y \\ E_z
\end{pmatrix} \longrightarrow
\vec{E}' = \begin{pmatrix}
E'_x \\ E'_y \\ E'_z
\end{pmatrix} = \begin{pmatrix}
E_x \\ \gamma \left( E_y - v B_z \right)\\ \gamma \left( E_z + v B_y \right)
\end{pmatrix}
\]
\[
\vec{B} = \begin{pmatrix}
B_x \\ B_y \\ B_z
\end{pmatrix} \longrightarrow
\vec{B}' = \begin{pmatrix}
B'_x \\ B'_y \\ B'_z
\end{pmatrix} = \begin{pmatrix}
B_x \\ \gamma \left( B_y + \frac{v}{c^2} E_z \right) \\ \gamma \left( B_z - \frac{v}{c^2} E_y \right)
\end{pmatrix}
\]

\noindent
Assume that in the laboratory frame $S$, there is an electric field $\vec{E} = E_0 \hat{y}$ and no magnetic field $\vec{B} = 0$. An inertial observer $S'$ moves relative to $S$ with constant velocity $v$ along the $x$-axis.
\begin{enumerate}
    \item[(a)] Use the field transformation expressions given above for $\vec{E}$ and $\vec{B}$ to find the electric and magnetic fields $\vec{E}'$ and $\vec{B}'$ as measured in the frame $S'$.
    \item[(b)] According to your result in (a), what does an observer in $S'$ measure for the electric field $\vec{E}'$ and the magnetic field $\vec{B}'$? Explain it.
\item[(c)] What is the force $\vec{F}'=q\left(\vec{E}' + \vec{u}'\times \vec{B}'\right)$ acting on the particle in the frame $S'$? Remember that the velocity $\vec{u}'$ of the particle at time $t'$ is given by the vector $\vec{u}'=\begin{pmatrix}
u'_x \\ u'_y \\ u'_z
\end{pmatrix}$.
    \item[(d)] (\textit{Hard-Optional}) At $t'=0$, a particle of charge $q$ and mass $m$ is at rest in the frame $S'$ at the origin, $\vec{u}'(0)=0$. Can you solve the equation of motions of the particle and obtain the velocity $\vec{u}'(t')$ of the particle in the frame $S'$? Explain in words why a purely electric field in one frame appears partly as a magnetic field in another, and how that “new” magnetic field influences the charged particle’s trajectory.

\end{enumerate}

\end{enumerate}



\end{document}
